% 2007 年湖南卷文科第 6 题：正四棱柱 ABCD-A_1B_1C_1D_1，E、F 分别为 AB_1、BC_1 的中点。
% 来源 PDF：普通高考/2007/2007湖南文.pdf，SHA256 8e889ae321ffa864b24c505b40e7e96956b301487f5a106264b7fcd9f943e068
% 原页：第 1 页（A3 三栏）第 1 栏下部，印刷页 381，图在题干下方。
% 重绘依据：斜投影正方体，A 左下、B 右下、C 右后、D 左后，上底同名加下标 1；
%   可见棱与上底对角线 A_1C_1 实线，被遮棱 D-D_1、C-D、D-A 与上底对角线 B_1D_1 虚线；
%   E 在 AB_1 上、F 在 BC_1 上，均取中点并画小黑点；标注位置按原图。
\begin{tikzpicture}[scale=1.05, semithick,
  every node/.style={inner sep=1.5pt}]
  \coordinate (A)  at (0,0);
  \coordinate (B)  at (2.4,0);
  \coordinate (D)  at (0.9,0.6);
  \coordinate (C)  at (3.3,0.6);
  \coordinate (A1) at (0,2.4);
  \coordinate (B1) at (2.4,2.4);
  \coordinate (D1) at (0.9,3.0);
  \coordinate (C1) at (3.3,3.0);
  % 被遮棱
  \draw[dashed] (D) -- (D1);
  \draw[dashed] (C) -- (D);
  \draw[dashed] (D) -- (A);
  % 可见棱
  \draw (A) -- (B) -- (C) -- (C1) -- (D1) -- (A1) -- cycle;
  \draw (A) -- (A1);
  \draw (B) -- (B1);
  \draw (C) -- (C1);
  \draw (B1) -- (C1);
  \draw (A1) -- (B1);
  % 上底对角线
  \draw (A1) -- (C1);
  \draw[dashed] (B1) -- (D1);
  % 两条侧面对角线（E、F 所在）
  \draw (A) -- (B1);
  \draw (B) -- (C1);
  % E、F
  \coordinate (E) at ($(A)!0.5!(B1)$);
  \coordinate (F) at ($(B)!0.5!(C1)$);
  \fill (E) circle (1.4pt);
  \fill (F) circle (1.4pt);
  \node[anchor=east]  at ([xshift=-2pt]E) {$E$};
  \node[anchor=west]  at ([xshift=2pt]F) {$F$};
  % 顶点标注
  \node[anchor=north east] at (A)  {$A$};
  \node[anchor=north west] at (B)  {$B$};
  \node[anchor=west]       at (C)  {$C$};
  \node[anchor=east]       at (D)  {$D$};
  \node[anchor=south east] at (A1) {$A_1$};
  \node[anchor=south west] at (B1) {$B_1$};
  \node[anchor=south west] at (C1) {$C_1$};
  \node[anchor=south east] at (D1) {$D_1$};
\end{tikzpicture}
